\section{Background and Related Work}
% Weakness of passwords
    % credential stuffing etc.

Passwords are notoriously weak as a sole authentication factor \cite{password_reuse, florencio_large_2006}, with attacks such as password spraying \cite{spraying} and credential stuffing \cite{credential_stuffing} affecting millions of online accounts each year. Nevertheless, a combination of circumstances has led passwords to remain the dominant authentication factor for online accounts today, making secure password management of paramount importance to one's overall security posture. An effective password management solution will allow users to have unique, high-entropy passwords for every website while maintaining security, portability, and ease of use.

In this section, we aim to motivate the need for a multi-factor deterministic password generator by discussing the range of existing password management and deterministic password generation solutions and their associated pitfalls.

\subsection{Password Managers}

The primary function of password managers is to provide an encrypted vault for the secure and portable storage of passwords. Today, users are faced with a difficult choice between the convenience of cloud-based password managers and the security of open-source self-managed solutions.

Cloud-based password managers, such as LastPass \cite{lastpass_arch}, Dashlane \cite{dashlane_arch}, and 1Password \cite{1password_arch}, are generally amongst the most popular password management solutions due to their perceived ease of use. The core technology used to secure these platforms is password-based key derivation.

Password-based key derivation functions like PBKDF2 \cite{rfc2898}, bcrypt \cite{bcrypt}, scrypt \cite{scrypt}, and Argon2 \cite{argon2} are one-way functions that convert a user's password into a cryptographic key that can be used for encryption.
Password-based key derivation functions are built upon cryptographic hash functions like SHA-256 \cite{rfc6234}, but also incorporate a degree of intentional computational inefficiency that increases the relative difficulty of brute-force attacks. For example, the PBKDF2 configuration used by LastPass uses 100,000 sequential rounds of SHA-256 to increase its computational difficulty.

Most cloud-based password managers use a password-based key derivation function such as PBKDF2 to derive a key from the user's password upon login.
A symmetric encryption function like AES-256 \cite{fips197} is then used to encrypt all of a user's secrets on the client side prior to their storage in a centralized database.
In theory, even an adversary with full access to the database will not be able to derive the key needed to decrypt the user's secrets without knowing their password.

Data breaches associated with major password managers are surprisingly common; LastPass alone has experienced at least 8 major security incidents \cite{lastpass_security}, including a recent total data breach of encrypted credentials \cite{toubba_notice_2022}. As no cloud-based service is totally immune from such a compromise, the security of these applications reduce, in practice, to the security of their users' master passwords. Given the aforementioned weakness of passwords as a sole authentication factor, there is a serious risk of attackers compromising the credentials stored in these services by performing offline brute-force attacks. As such, cloud-based password managers may even constitute a net liability for some users by creating a central point of failure.

Self-hosted open-source password managers like Bitwarden \cite{bitwarden}, KeePass \cite{reichl_keepass_nodate}, and KeePassX \cite{keepassx} use a similar architecture to cloud-based password managers to encrypt and store credentials using password-derived keys. However, these tools aim to address the vulnerability of cloud-based password managers to massive centralized data breaches by only storing encrypted passwords locally on a user's device, or on a server directly owned and maintained by the end user. 
In doing so, they avoid creating a concentrated high-value attack target as is inevitably the case with cloud-based password managers.

Still, self-hosted password managers are not immune to attack, with threats such as malware posing a risk to the database of encrypted credentials.
As with cloud-based password managers, the user's security in the event of a breach ultimately reduces to that of their master password.

Moreover, open-source password management solutions have failed to achieve widespread adoption due their perceived impracticality. The lack of a centralized database to store credentials requires users to either host and maintain their own servers, or to manually synchronize data between all of their devices. Further, while cloud-based password managers are usually designed around resilient and highly-available architectures, self-hosted solutions can be susceptible to a total loss of data if a single device is lost or rendered inoperable.
As such, most users have gravitated toward centralized password managers despite their relative security drawbacks.

\subsection{Deterministic Password Generators}
\label{sec:survey}

Deterministic password generators (DPGs) represent an interesting alternative to both cloud-based and self-hosted password management applications. Instead of storing encrypted passwords in any vault, DPGs work by deriving site-specific passwords using a cryptographic hash of the site's domain name and the user's master password. The deterministic nature of the underlying hash function ensures that the DPG will always generate the same password for a given site as long as the user remembers their master password, without needing to synchronize additional values between devices.
This basic idea was originally proposed by HP in 2002 \cite{karp2002site}, and was popularized by Stanford's PwdHash in 2003 \cite{PwdHash}. When implemented as a browser extension that automatically determines a site's URL, it has the additional benefit of resisting phishing attacks.

Since the advent of DPGs in the early 2000s, dozens of DPG tools have been developed and released as open-source websites, mobile applications, and browser extensions. Despite this, DPGs have seen even less mainstream adoption than self-hosted password managers, with the 20 most popular DPG extensions having less than 10,000 combined downloads.

Consumers' unwillingness to adopt DPGs may stem from a variety of security, privacy, and usability flaws \cite{noauthor_4_2016, palant_security_2016} that researchers have discovered with open-source DPGs.
However, the relative obscurity of DPGs as a mainstream tool has resulted in limited literature being generated on this topic.
To gain a better understanding of the current landscape of DPGs and their associated flaws, we performed a large-scale analysis of existing DPG implementations.
By searching GitHub, the Chrome Web Store, and the Firefox Extensions website, we identified 45 free software packages implementing some form of deterministic password generation algorithm.

We explored the source code of these 45 DPG extensions to identify whether there were any particular threats to security, privacy, or usability that may limit their use as an alternative to traditional password managers. Our findings for each DPG system are detailed in Table \ref{tab:dpg-survey}. Overall, we found four key issues affecting a large portion of the studied implementations:

\begin{enumerate}[leftmargin=*]
    \itemsep 1em
    \item \textbf{Brute-Force Susceptibility}. By far the largest concern with current DPG implementations is the ability to attack a user's master password. Because of the lack of stored ciphertexts in DPGs, exposure of the master password is sufficient to reveal all of a user's site passwords.    
    However, as it stands, obtaining any of a user's site-specific passwords allows an attacker to perform an offline brute-force attack on that user's master password by checking which master password would, when combined with the site's known URL, have resulted in the correct site password being generated.
    
    \medskip
    
    This threat would be somewhat mitigated by the use of a strong progressive password-based key derivation function to increase the computational difficulty of a brute-force attack. Unfortunately, of the 45 applications surveyed, 28 only used a standard cryptographic hash, with 11 using deprecated function (MD5 or SHA1). A further 13 applications did use a progressive hash function, but supplied a cost parameter too low to effectively deter modern hardware. Only four of the 45 applications implemented a progressive hash function with well-chosen cost parameters.
    
    \item \textbf{No Policy Support}. A second issue that affects most DPGs is the lack of support for applications with complex password policies.
    Websites often enforce a number of password strength and complexity requirements, with length, capitalization, character sets, etc. varying greatly from one service to another. No single static generator can hope to accomodate a wide range password policies without tuning the generation algorithm for each site.
    
    \medskip
    
    Of the 45 DPGs surveyed, 27 provide no settings at all for customizing password generation; users of these applications would likely face scenarios in which the DPG cannot generate a password compatible with a service they wish to use. The remaining 18 provide limited customization, which in most cases only allows the output length to be tweaked. Users of these applications must also remember these settings and configure them correctly on every login.
    
    \item \textbf{No Revocation Support}. Similarly, websites often impose password rotation policies that require users to periodically change their password. However, the deterministic nature of DPGs usually does not allow users to have more than one password for a given website. Only four of the DPGs we analyzed provided a way to revoke a password and generate a new one for the same service. Of these, two required users to remember a revocation counter value for each website, while the other two stored this counter value in a file. 
    
    \item \textbf{Multi-Factor Authentication}. Finally, a fundamental limitation of all existing DPGs is the solitary reliance on a master password and the lack of support for multi-factor authentication. Most centralized password managers are enhanced by their support for multi-factor authentication at the login stage, though password vaults are still ultimately only encrypted with password-derived keys.
\end{enumerate}

In addition to the common issues noted above, we found specific vulnerabilities in a six implementations: two schemes designed their own insecure hash function, two incorrectly deployed progressive hash functions, and two had obvious network-related vulnerabilities. More information about these issues is given in the footnotes of Table \ref{tab:dpg-survey}.

As a result of many of the discussed issues, DPGs have remained a relatively obscure technology and have largely been dismissed as a viable alternative to conventional password managers. Thus, DPG projects have largely been abandoned, with most open-source DPGs receiving no updates in the past five years, as illustrated in Table \ref{tab:dpg-survey}. We hope, in this paper, to revive the field of deterministic password generation, aided by recent cryptographic advances that enable the incorporation of multiple authentication factors into derived keys.

\subsection{Multi-Factor Authentication}
% , including TOTP, HOTP, OOBA (e.g., Email/SMS), and HMAC-SHA1 (e.g., YubiKey)

The most popular form of multi-factor authentication in use today is ``out-of-band authentication'' (OOBA) \cite{ooba}, which includes email and SMS-based authentication. However, these channels are not trustless and fundamentally require a server; thus, they do not receive significant attention in this paper.

Aside from OOBA, popular MFA options include ``soft tokens'' like HMAC-based one-time password (HOTP) \cite{rfc4226} and time-based one-time password (TOTP) \cite{rfc6238}, as implemented in applications like Google Authenticator. These factors use a counter or timestamp to generate one-time passwords based on a pre-shared secret.
``Hard tokens'' such as YubiKeys \cite{yubikey} are also a popular option that requires using specialized hardware.

\subsection{Multi-Factor Key Derivation}
% % MFKDF
The Multi-Factor Key Derivation Function (MFKDF) \cite{mfkdf} is a recent improvement over PBKDFs that incorporates multiple authentication factors into the key derivation process. It is a trustless cryptographic operation that can handle many popular authentication factors like HOTP, TOTP, and YubiKeys on the client side without the need for a trusted server.

The MFKDF specification contains two major architectural components. The first component is the set of so-called ``factor constructions,'' which convert a dynamic \textit{factor witness}\footnote{In this case, the \textit{witness} refers to the message used to authenticate (e.g., a 6-digit OTP), which is often not the same as the underlying shared secret.} and public parameters into static key material. The public parameters require no security assumptions and can safely be stored in a database without concern for revealing information about the factors to potential adversaries. Constructions are given for a variety of popular authentication factors.

The second major component of MFKDF is the key derivation function itself, which adds a secret sharing layer to provide functionality such as threshold-based key derivation, advanced policy enforcement, and factor recovery. Together, these components convert multiple authentication factors into a cryptographic key, and serve as a replacement for PBKDFs. 

This paper proposes using MFKDF to enhance the security and utility of classical DPGs. While DPGs have been around in some form since 2002, we believe the recent introduction of MFKDF has provided a critical tool for their practical use.

\eject